\documentclass[addpoints]{physlabs}

\labtitle{Voltage \& Electrostatic Potential}
\labnum{7}
\labnumroman{VII}
\course{PHYSICS 122 - 142 - 182}
\coursesubtitle{Electricity and Magnetism}
\date{\today}

\begin{document}

\maketitle

% PRELAB SECTION
\begin{prelab}

    \filbreak
    \question[1] In the first part of this lab, you will be using a mechanically-balanced, axially-symmetric cylindrical capacitor to measure an absolute voltage. Each time you take a measurement, but before you read the microscope scale, you should perform a critical procedural technique.  What is it, why is it critical, and how many times should you perform it?
    \begin{solution}[5cm]
        
    \end{solution}

    \filbreak
    \question[1] How do you determine the electric field lines for the second part of the experiment? Electric fields are vector quantities.  How will you determine the direction that they should point?
    \begin{solution}[5cm]
        
    \end{solution}

\end{prelab}

% LAB SECTION
\begin{labsection}

\begin{objective}
    The purpose is to calculate an absolute voltage measurement by mechanical means and illustrate the concepts of the electric field by means of experimental demonstration.
\end{objective}
    
    \begin{theory}
    A particle of charge $q$ moving with a velocity $\boldsymbol{v}$ through a region of space containing both electric, $\boldsymbol{E}$, and magnetic, $\boldsymbol{B}$, fields experiences the Lorentz force
    \begin{align}
        \boldsymbol{F} = q(\boldsymbol{E}+\boldsymbol{v}\times\boldsymbol{B}).
    \end{align}
    Our experimental apparatus first accelerates electrons ($q = -e=-1.60217663\times 10^{-19} C$) through a region containing only an electric field.
    The electric field tunes their velocity, so that when they enter another region containing only a magnetic field they move in a circle whose size depends on the charge to mass ratio, $e/m$.
    Let's understand why.
    
    A cathode within the apparatus heats up, causing electrons to be emitted. 
    The cathode is connected to the negative of the high voltage supply, $V_-$. 
    The anode is connected to the high voltage positive end of the supply, $V_+$, a distance $d$ away from the cathode.
    The potential difference sets up an electric field of strength
    \begin{align}
        E = \frac{V_--V_+}{d}=-\frac{\Delta V}{d}.
    \end{align}
    The minus sign is because the field points to the left (see figure below), from positive to negative.
    Electrons are rapidly accelerated toward the anode by the electric field. 
    Most of the electrons hit the anode, but some of them pass through a hole in the anode and form the electron beam.
        \begin{figure}[h]
        \centering
        \includegraphics[width=0.3\linewidth]{Lab 07: Electrostatic Potential/CathodeAnodeElectronBeam.png}
        \label{fig:placeholder}
    \end{figure}
    
    By Newton's second law, the acceleration they experience due to the electric field is
    \begin{align}
        a = \frac{q}{m}E=\frac{e\Delta V}{md},
    \end{align}
    which, between the two plates, is constant.
    The electrons emitted by the cathode move with some initial velocity $v_i$.
    We can use kinematics to find the final speed of the elctrons when they pass through the hole in the anode. 
    Plugging into the kinematic equation
    \begin{align}
        v_f^2 = v_i^2+2ad=v_i^2+\frac{2e\Delta V}{m}
    \end{align}
    The final speed of the electrons is much greater than what they started with, $v_f\gg v_i$, so we can approximate their velocity leaving the anode as
    \begin{align}
        v_f \approx \sqrt{\frac{2e\Delta V}{m}}
    \end{align}
    

    
    A pair of Helmholtz coils provides the uniform magnetic field oriented perpendicular to the path of the electron. 
    There is no electric field present. 
    For this case, the Lorentz force reduces to
    \end{theory}
    
    \begin{equipment}
    \item lalala
    \end{equipment}
    
    \begin{experiment}
    \end{experiment}
\end{labsection}

\begin{postlab}

    \filbreak
    \question[1] In Graph~\ref{graph: two point field lines}, sketch the electric field lines according to your measured mapping of equipotential surfaces.  Be sure to include arrows to show the direction of your field lines.

    \begin{center}
        \captionsetup{font=normalsize}
        \captionof{graph}{Two-point potential and electric lines. \label{graph: two point field lines}}
        \includegraphics[width=\linewidth]{Lab 07: Electrostatic Potential/graphite-paper-grid.png}
    \end{center}

    \filbreak
    \question[1] In Graph~\ref{graph: parallel field lines}, sketch the electric field lines according to your measured mapping of equipotential surfaces.  Be sure to include arrows to show the direction of your field lines.

    \begin{center}
        \captionsetup{font=normalsize}
        \captionof{graph}{Parallel-plate potential and electric lines. \label{graph: parallel field lines}}
        \includegraphics[width=\linewidth]{Lab 07: Electrostatic Potential/graphite-paper-grid.png}
    \end{center}

\end{postlab}

\end{document}